\section{信頼性と故障}

\subsection{平成30年度 問2-2}
\label{q:h30-reg-basic-2-2}

\noindent\textbf{原典：} sources/H30/basic.pdf、PDF物理7頁。

\textbf{公式正答：} 3（single）。

\par\textit{原典本文の校正と独自解説は未作成。分類のみの登録項目。}

\section{工程管理・PERT}

\subsection{平成26年度 問2-1}
\label{q:h26-reg-basic-2-1}

\noindent\textbf{原典：} sources/H26/basic.pdf、PDF物理5頁。

\textbf{公式正答：} 1（single）。

\par\textit{原典本文の校正と独自解説は未作成。分類のみの登録項目。}

\section{暗号・認証・攻撃対策}

\subsection{平成23年度 問2-5}
\label{q:h23-reg-basic-2-5}

\noindent\textbf{原典：} sources/H23/basic.pdf、PDF物理7頁。

\textbf{公式正答：} 2（single）。

\par\textit{原典本文の校正と独自解説は未作成。分類のみの登録項目。}

\subsection{平成27年度 問2-5}
\label{q:h27-reg-basic-2-5}

\noindent\textbf{原典：} sources/H27/basic.pdf、PDF物理6頁。

\textbf{公式正答：} 2（single）。

\par\textit{原典本文の校正と独自解説は未作成。分類のみの登録項目。}

\subsection{平成29年度 問2-1}
\label{q:h29-reg-basic-2-1}

\noindent\textbf{原典：} sources/H29/basic.pdf、PDF物理6頁。

\textbf{公式正答：} 1（single）。

\par\textit{原典本文の校正と独自解説は未作成。分類のみの登録項目。}

\subsection{平成30年度 問2-1}
\label{q:h30-reg-basic-2-1}

\noindent\textbf{原典：} sources/H30/basic.pdf、PDF物理6頁。

\textbf{公式正答：} 4（single）。

\par\textit{原典本文の校正と独自解説は未作成。分類のみの登録項目。}

\subsection{令和元年度再試験 問2-1}
\label{q:r01-ret-basic-2-1}

\noindent\textbf{原典：} sources/R01\_retest/basic.pdf、PDF物理6頁。

\textbf{公式正答：} 5（single）。

\par\textit{原典本文の校正と独自解説は未作成。分類のみの登録項目。}

\subsection{令和2年度 問2-3}
\label{q:r02-reg-basic-2-3}

\noindent\textbf{原典：} sources/R02/basic.pdf、PDF物理8頁。

\textbf{公式正答：} 1（single）。

\par\textit{原典本文の校正と独自解説は未作成。分類のみの登録項目。}

\subsection{令和3年度 問2-1}
\label{q:r03-reg-basic-2-1}

\noindent\textbf{原典：} sources/R03/basic.pdf、PDF物理6頁。

\textbf{公式正答：} 1（single）。

\par\textit{原典本文の校正と独自解説は未作成。分類のみの登録項目。}

\subsection{令和6年度 問2-6}
\label{q:r06-reg-basic-2-6}

\noindent\textbf{原典：} sources/R06/basic.pdf、PDF物理9頁。

\textbf{公式正答：} 4（single）。

\par\textit{原典本文の校正と独自解説は未作成。分類のみの登録項目。}

\subsection{令和7年度 問2-4}
\label{q:r07-reg-basic-2-4}

\noindent\textbf{原典：} sources/R07/basic.pdf、PDF物理6頁。

\textbf{公式正答：} 5（single）。

\par\textit{原典本文の校正と独自解説は未作成。分類のみの登録項目。}

\section{IP・通信}

\subsection{平成23年度 問2-1}
\label{q:h23-reg-basic-2-1}

\noindent\textbf{原典：} sources/H23/basic.pdf、PDF物理5頁。

\textbf{公式正答：} 4（single）。

\par\textit{原典本文の校正と独自解説は未作成。分類のみの登録項目。}

\subsection{平成23年度 問2-2}
\label{q:h23-reg-basic-2-2}

\noindent\textbf{原典：} sources/H23/basic.pdf、PDF物理5頁。

\textbf{公式正答：} 2（single）。

\par\textit{原典本文の校正と独自解説は未作成。分類のみの登録項目。}

\subsection{平成26年度 問2-6}
\label{q:h26-reg-basic-2-6}

\noindent\textbf{原典：} sources/H26/basic.pdf、PDF物理8頁。

\textbf{公式正答：} 3（single）。

\par\textit{原典本文の校正と独自解説は未作成。分類のみの登録項目。}

\subsection{平成28年度 問2-6}
\label{q:h28-reg-basic-2-6}

\noindent\textbf{原典：} sources/H28/basic.pdf、PDF物理11頁。

\textbf{公式正答：} 5（single）。

\par\textit{原典本文の校正と独自解説は未作成。分類のみの登録項目。}

\subsection{令和元年度 問2-5}
\label{q:r01-reg-basic-2-5}

\noindent\textbf{原典：} sources/R01/basic.pdf、PDF物理10頁。

\textbf{公式正答：} 5（single）。

\par\textit{原典本文の校正と独自解説は未作成。分類のみの登録項目。}

\subsection{令和3年度 問2-3}
\label{q:r03-reg-basic-2-3}

\noindent\textbf{原典：} sources/R03/basic.pdf、PDF物理7頁。

\textbf{公式正答：} 2（single）。

\par\textit{原典本文の校正と独自解説は未作成。分類のみの登録項目。}

\subsection{令和4年度 問2-4}
\label{q:r04-reg-basic-2-4}

\noindent\textbf{原典：} sources/R04/basic.pdf、PDF物理7頁。

\textbf{公式正答：} 5（single）。

\par\textit{原典本文の校正と独自解説は未作成。分類のみの登録項目。}

\subsection{令和4年度 問2-6}
\label{q:r04-reg-basic-2-6}

\noindent\textbf{原典：} sources/R04/basic.pdf、PDF物理9頁。

\textbf{公式正答：} 4（single）。

\par\textit{原典本文の校正と独自解説は未作成。分類のみの登録項目。}

\subsection{令和6年度 問2-3}
\label{q:r06-reg-basic-2-3}

\noindent\textbf{原典：} sources/R06/basic.pdf、PDF物理7頁。

\textbf{公式正答：} 2（single）。

\par\textit{原典本文の校正と独自解説は未作成。分類のみの登録項目。}

\subsection{令和7年度 問2-3}
\label{q:r07-reg-basic-2-3}

\noindent\textbf{原典：} sources/R07/basic.pdf、PDF物理6頁。

\textbf{公式正答：} 3（single）。

\par\textit{原典本文の校正と独自解説は未作成。分類のみの登録項目。}

\subsection{令和7年度 問2-5}
\label{q:r07-reg-basic-2-5}

\noindent\textbf{原典：} sources/R07/basic.pdf、PDF物理7頁。

\textbf{公式正答：} 4（single）。

\par\textit{原典本文の校正と独自解説は未作成。分類のみの登録項目。}

\section{論理式・命題}

\subsection{平成23年度 問2-3}
\label{q:h23-reg-basic-2-3}

\noindent\textbf{原典：} sources/H23/basic.pdf、PDF物理6頁。

\textbf{公式正答：} 2（single）。

\par\textit{原典本文の校正と独自解説は未作成。分類のみの登録項目。}

\subsection{平成25年度 問2-3}
\label{q:h25-reg-basic-2-3}

\noindent\textbf{原典：} sources/H25/basic.pdf、PDF物理8頁。

\textbf{公式正答：} 1（single）。

\par\textit{原典本文の校正と独自解説は未作成。分類のみの登録項目。}

\subsection{平成25年度 問2-4}
\label{q:h25-reg-basic-2-4}

\noindent\textbf{原典：} sources/H25/basic.pdf、PDF物理8頁。

\textbf{公式正答：} 4（single）。

\par\textit{原典本文の校正と独自解説は未作成。分類のみの登録項目。}

\subsection{平成28年度 問2-2}
\label{q:h28-reg-basic-2-2}

\noindent\textbf{原典：} sources/H28/basic.pdf、PDF物理8頁。

\textbf{公式正答：} 2（single）。

\par\textit{原典本文の校正と独自解説は未作成。分類のみの登録項目。}

\subsection{平成29年度 問2-4}
\label{q:h29-reg-basic-2-4}

\noindent\textbf{原典：} sources/H29/basic.pdf、PDF物理9頁。

\textbf{公式正答：} 3（single）。

\par\textit{原典本文の校正と独自解説は未作成。分類のみの登録項目。}

\subsection{平成30年度 問2-4}
\label{q:h30-reg-basic-2-4}

\noindent\textbf{原典：} sources/H30/basic.pdf、PDF物理8頁。

\textbf{公式正答：} 5（single）。

\par\textit{原典本文の校正と独自解説は未作成。分類のみの登録項目。}

\subsection{令和元年度再試験 問2-5}
\label{q:r01-ret-basic-2-5}

\noindent\textbf{原典：} sources/R01\_retest/basic.pdf、PDF物理10頁。

\textbf{公式正答：} 1（single）。

\par\textit{原典本文の校正と独自解説は未作成。分類のみの登録項目。}

\subsection{令和2年度 問2-2}
\label{q:r02-reg-basic-2-2}

\noindent\textbf{原典：} sources/R02/basic.pdf、PDF物理7頁。

\textbf{公式正答：} 3（single）。

\par\textit{原典本文の校正と独自解説は未作成。分類のみの登録項目。}

\subsection{令和3年度 問2-2}
\label{q:r03-reg-basic-2-2}

\noindent\textbf{原典：} sources/R03/basic.pdf、PDF物理6頁。

\textbf{公式正答：} 2（single）。

\par\textit{原典本文の校正と独自解説は未作成。分類のみの登録項目。}

\subsection{令和3年度 問2-4}
\label{q:r03-reg-basic-2-4}

\noindent\textbf{原典：} sources/R03/basic.pdf、PDF物理8頁。

\textbf{公式正答：} 2（single）。

\par\textit{原典本文の校正と独自解説は未作成。分類のみの登録項目。}

\subsection{令和5年度 問2-5}
\label{q:r05-reg-basic-2-5}

\noindent\textbf{原典：} sources/R05/basic.pdf、PDF物理8頁。

\textbf{公式正答：} 1（single）。

\par\textit{原典本文の校正と独自解説は未作成。分類のみの登録項目。}

\section{キャッシュ・CPU}

\subsection{平成23年度 問2-4}
\label{q:h23-reg-basic-2-4}

\noindent\textbf{原典：} sources/H23/basic.pdf、PDF物理6頁。

\textbf{公式正答：} 1（single）。

\par\textit{原典本文の校正と独自解説は未作成。分類のみの登録項目。}

\subsection{平成25年度 問2-6}
\label{q:h25-reg-basic-2-6}

\noindent\textbf{原典：} sources/H25/basic.pdf、PDF物理9頁。

\textbf{公式正答：} 3（single）。

\par\textit{原典本文の校正と独自解説は未作成。分類のみの登録項目。}

\subsection{平成28年度 問2-4}
\label{q:h28-reg-basic-2-4}

\noindent\textbf{原典：} sources/H28/basic.pdf、PDF物理10頁。

\textbf{公式正答：} 4（single）。

\par\textit{原典本文の校正と独自解説は未作成。分類のみの登録項目。}

\subsection{平成29年度 問2-6}
\label{q:h29-reg-basic-2-6}

\noindent\textbf{原典：} sources/H29/basic.pdf、PDF物理10頁。

\textbf{公式正答：} 2（single）。

\par\textit{原典本文の校正と独自解説は未作成。分類のみの登録項目。}

\subsection{令和2年度 問2-6}
\label{q:r02-reg-basic-2-6}

\noindent\textbf{原典：} sources/R02/basic.pdf、PDF物理10頁。

\textbf{公式正答：} 5（single）。

\par\textit{原典本文の校正と独自解説は未作成。分類のみの登録項目。}

\section{数値表現}

\subsection{平成24年度 問2-1}
\label{q:h24-reg-basic-2-1}

\noindent\textbf{原典：} sources/H24/basic.pdf、PDF物理5頁。

\textbf{公式正答：} 4（single）。

\paragraph{問題の要旨} 数値表現と誤差に関する5記述から、最も適切なものを選ぶ。

\paragraph{問われる核心} 有限精度計算で誤差を生む原因

\paragraph{必要知識} 有効桁数、丸め、桁落ち、情報落ち、指数範囲。

\paragraph{考え方（査読待ち初稿）} 有限精度は固定・浮動のどちらでも誤差を生じ得る。浮動小数点には有限の指数範囲もある。4の「改善できることが多い」は仮数精度を増す効果として適切である。5は大小の加算で小さい量が失われる例である。

\paragraph{正答を選ぶ理由} 有限精度は固定・浮動のどちらでも誤差を生じ得る。浮動小数点には有限の指数範囲もある。4の「改善できることが多い」は仮数精度を増す効果として適切である。5は大小の加算で小さい量が失われる例である。

\paragraph{選択肢の検討（査読待ち初稿）}

\noindent\textbf{1：不適切な解答候補} 浮動小数点では桁あふれ・桁不足はない。指数の範囲は有限であり、上限・下限の範囲超過がある。\par

\noindent\textbf{2：不適切な解答候補} 桁落ちは乗算又は除算で生じる。近い数の減算又は逆符号の加算で有効桁が相殺される現象である。\par

\noindent\textbf{3：不適切な解答候補} 固定小数点なら誤差はない。小数桁数が有限なら表現できない値があり、演算後の丸めもある。\par

\noindent\textbf{4：適切な解答候補} 単精度の丸め誤差は倍精度により改善できることが多い。通常は仮数の桁数が増え、同じ計算の表現・丸め誤差を小さくできる。ただし常に完全に解消するという主張ではない。\par

\noindent\textbf{5：不適切な解答候補} 10進7桁なら100000.0+0.01に誤差はない。正確な100000.01には8桁の有効数字が必要で、7桁では0.01が保持されない。\par

\paragraph{再利用する考え方} 表現・演算の丸め、相殺による桁落ち、大小の加算による情報落ち、範囲超過を区別する。

\paragraph{共通理論} \hyperref[theory:information:error]{第
ef*{theory:information:error}節}

\par 参考：\cite{goldberg1991}。

\par 難易度 2（校準初稿）、優先度未評価。

\subsection{平成25年度 問2-5}
\label{q:h25-reg-basic-2-5}

\noindent\textbf{原典：} sources/H25/basic.pdf、PDF物理9頁。

\textbf{公式正答：} 2（single）。

\paragraph{問題の要旨} 0.85を二進小数4桁で切り捨て、その値を半分にして再度4桁で切り捨てる。

\paragraph{問われる核心} 有限二進小数の量子化と演算順

\paragraph{必要知識} 二進小数の位取り、切捨てによる量子化。

\paragraph{考え方（査読待ち初稿）} 初回は \(\lfloor16\times0.85\rfloor/16=13/16=(0.1101)_2\)。これを半分にすると13/32=(0.01101)₂である。4桁に切り捨てると6/16=(0.0110)₂=0.375になる。

\paragraph{正答を選ぶ理由} 初回は \(\lfloor16\times0.85\rfloor/16=13/16=(0.1101)_2\)。これを半分にすると13/32=(0.01101)₂である。4桁に切り捨てると6/16=(0.0110)₂=0.375になる。

\paragraph{選択肢の検討（査読待ち初稿）}

\noindent\textbf{1：不適切な解答候補} 0.1101／0.0101／0.375。除算後の0110を0101としている。0101の値も0.3125であり提示の0.375に合わない。\par

\noindent\textbf{2：適切な解答候補} 0.1101／0.0110／0.375。13/16を半分にし、6/16へ切り捨てる両段階と一致する。\par

\noindent\textbf{3：不適切な解答候補} 0.1101／0.0110／0.425。0.425は元の0.85を正確に半分にした値であり、段階ごとの切捨てを無視している。\par

\noindent\textbf{4：不適切な解答候補} 0.1010／0.0101／0.375。初回は1101であり1010ではない。除算後の0101も誤りである。\par

\noindent\textbf{5：不適切な解答候補} 0.1010／0.0101／0.425。初回と除算後のビットが誤りで、0.425も切捨て後の値ではない。\par

\paragraph{再利用する考え方} 位取りの基数と、小数の有限表現が可能かどうかを区別する。小数は2倍操作で上位桁から求める。

\paragraph{共通理論} \hyperref[theory:information:radix]{第
ef*{theory:information:radix}節}

\par 難易度 2（校準初稿）、優先度未評価。

\subsection{平成27年度 問2-1}
\label{q:h27-reg-basic-2-1}

\noindent\textbf{原典：} sources/H27/basic.pdf、PDF物理5頁。

\textbf{公式正答：} 5（single）。

\paragraph{問題の要旨} 二進計算機で0.01を100回加え、正しいプログラムでも結果が1と一致しない原因を選ぶ。

\paragraph{問われる核心} 進数依存の表現誤差と累積

\paragraph{必要知識} 有限二進小数の条件、入力と演算の丸め。

\paragraph{考え方（査読待ち初稿）} 0.01=1/100の既約分母には5が残り、二進有限小数にはならない。格納値は近似であり、加算の結果も有限桁へ丸められるので、100回の累積値が数学的な1と一致しないことがある。誤差の符号や実際の値は形式・演算順に依存する。

\paragraph{正答を選ぶ理由} 0.01=1/100の既約分母には5が残り、二進有限小数にはならない。格納値は近似であり、加算の結果も有限桁へ丸められるので、100回の累積値が数学的な1と一致しないことがある。誤差の符号や実際の値は形式・演算順に依存する。

\paragraph{選択肢の検討（査読待ち初稿）}

\noindent\textbf{1：不適切な解答候補} 桁落ち。ほぼ等しい量を減算して有効桁が失われる状況ではない。\par

\noindent\textbf{2：不適切な解答候補} 情報落ち。著しく桁の異なる量を加算して小さい量が消える原因を主題としていない。\par

\noindent\textbf{3：不適切な解答候補} オーバーフロー。1付近の合計が通常の浮動小数点の上限を超える状況ではない。\par

\noindent\textbf{4：不適切な解答候補} アンダーフロー。0.01は通常の浮動小数点の最小正数より十分大きい。\par

\noindent\textbf{5：適切な解答候補} 丸め誤差。0.01は二進有限小数で表せず、入力表現と各加算の丸めが蓄積し得る。\par

\paragraph{再利用する考え方} 表現・演算の丸め、相殺による桁落ち、大小の加算による情報落ち、範囲超過を区別する。

\paragraph{共通理論} \hyperref[theory:information:error]{第
ef*{theory:information:error}節}

\par 参考：\cite{goldberg1991}。

\par 難易度 2（校準初稿）、優先度未評価。

\subsection{平成27年度 問2-3}
\label{q:h27-reg-basic-2-3}

\noindent\textbf{原典：} sources/H27/basic.pdf、PDF物理6頁。

\textbf{公式正答：} 3（single）。

\paragraph{問題の要旨} 10進0.1を二進小数で小数点以下16桁まで表す。

\paragraph{問われる核心} 循環二進小数の生成

\paragraph{必要知識} 小数の反復2倍法、分母と有限展開。

\paragraph{考え方（査読待ち初稿）} 小数を順に2倍し整数部を取り出す。0.1→0.2→0.4→0.8→1.6で最初の4桁は0001である。残り0.6から1.2、0.4、0.8、1.6と循環し、0011の反復が現れる。16桁は0001100110011001となる。

\paragraph{正答を選ぶ理由} 小数を順に2倍し整数部を取り出す。0.1→0.2→0.4→0.8→1.6で最初の4桁は0001である。残り0.6から1.2、0.4、0.8、1.6と循環し、0011の反復が現れる。16桁は0001100110011001となる。

\paragraph{選択肢の検討（査読待ち初稿）}

\noindent\textbf{1：不適切な解答候補} 0.0000011001100110。先頭の1が6桁目にあり、0.1の4桁目と合わない。\par

\noindent\textbf{2：不適切な解答候補} 0.0000110011001100。先頭の1が5桁目で、0.1の表現を一桁右へずらした形である。\par

\noindent\textbf{3：適切な解答候補} 0.0001100110011001。反復2倍の出力0001100110011001に一致する。\par

\noindent\textbf{4：不適切な解答候補} 0.0011001100110011。先頭の1が3桁目で、0.1の表現を一桁左へずらした形である。\par

\noindent\textbf{5：不適切な解答候補} 0.0110011001100110。先頭の1が2桁目で、0.1より大きい値である。\par

\paragraph{再利用する考え方} 位取りの基数と、小数の有限表現が可能かどうかを区別する。小数は2倍操作で上位桁から求める。

\paragraph{共通理論} \hyperref[theory:information:radix]{第
ef*{theory:information:radix}節}

\par 難易度 2（校準初稿）、優先度未評価。

\subsection{平成27年度 問2-4}
\label{q:h27-reg-basic-2-4}

\noindent\textbf{原典：} sources/H27/basic.pdf、PDF物理6頁。

\textbf{公式正答：} 1（single）。

\paragraph{問題の要旨} 10進表記で2 TBの記憶容量を、問題が定義する2進表記のTBへ換算する。

\paragraph{問われる核心} 情報容量の接頭語と単位換算

\paragraph{必要知識} TBとTiBの単位量、情報量保存。

\paragraph{考え方（査読待ち初稿）} 容量をバイトへ統一してから新単位で割る。\(2\times10^{12}/2^{40}=2(1000/1024)^4\simeq1.819\)であり、約1.8となる。原典の「2進のTB」を現在の用語で明確に区別すればTiBである。

\paragraph{正答を選ぶ理由} 容量をバイトへ統一してから新単位で割る。\(2\times10^{12}/2^{40}=2(1000/1024)^4\simeq1.819\)であり、約1.8となる。原典の「2進のTB」を現在の用語で明確に区別すればTiBである。

\paragraph{選択肢の検討（査読待ち初稿）}

\noindent\textbf{1：適切な解答候補} 1.8 TB。2×10\textasciicircum{}12/2\textasciicircum{}40≈1.819を小数1桁へ概算した値に一致する。\par

\noindent\textbf{2：不適切な解答候補} 2.0 TB。2.0 TBは2×10\textasciicircum{}12/2\textasciicircum{}40≈1.819の概算と合わない。単位当たりの容量が大きくなるので数値は2より小さくなる。\par

\noindent\textbf{3：不適切な解答候補} 2.1 TB。2.1 TBは2×10\textasciicircum{}12/2\textasciicircum{}40≈1.819の概算と合わない。単位当たりの容量が大きくなるので数値は2より小さくなる。\par

\noindent\textbf{4：不適切な解答候補} 2.2 TB。2.2 TBは2×10\textasciicircum{}12/2\textasciicircum{}40≈1.819の概算と合わない。単位当たりの容量が大きくなるので数値は2より小さくなる。\par

\noindent\textbf{5：不適切な解答候補} 2.3 TB。2.3 TBは2×10\textasciicircum{}12/2\textasciicircum{}40≈1.819の概算と合わない。単位当たりの容量が大きくなるので数値は2より小さくなる。\par

\paragraph{再利用する考え方} 単位あたりのバイト数を先に定義し、容量の数値だけを比べない。

\paragraph{共通理論} \hyperref[theory:information:capacity]{第
ef*{theory:information:capacity}節}

\par 参考：\cite{nist_binary_prefix}。

\par 難易度 2（校準初稿）、優先度未評価。

\subsection{平成28年度 問2-3}
\label{q:h28-reg-basic-2-3}

\noindent\textbf{原典：} sources/H28/basic.pdf、PDF物理9頁。

\textbf{公式正答：} 1（single）。

\paragraph{問題の要旨} 11を例に、2で割った余りを配列へ格納する二進変換流れ図の分岐と更新を選ぶ。

\paragraph{問われる核心} 整数商と余りによる基数変換

\paragraph{必要知識} 除法の等式N=2q+r、0≤r<2、終端条件、配列の順序。

\paragraph{考え方（査読待ち初稿）} Nを11→5→2→1→0と整数商で更新し、余りは1、1、0、1の順で得る。格納順は下位桁からなので逆順表示により1011となる。正のNで継続し、N=0で停止する。N=0を入力として許す実装ではゼロ表示の別処理が必要だが、例示の正の入力の空欄判定には影響しない。

\paragraph{正答を選ぶ理由} Nを11→5→2→1→0と整数商で更新し、余りは1、1、0、1の順で得る。格納順は下位桁からなので逆順表示により1011となる。正のNで継続し、N=0で停止する。N=0を入力として許す実装ではゼロ表示の別処理が必要だが、例示の正の入力の空欄判定には影響しない。

\paragraph{選択肢の検討（査読待ち初稿）}

\noindent\textbf{1：適切な解答候補} N>0／N=0／Nを2で割った商。正のNについて余りを保存し整数商へ更新、商0で終了する。\par

\noindent\textbf{2：不適切な解答候補} N>0／N=0／Nを2で割った余り。余りをNへ戻すと桁の取り出しが進まず、余り1では停止しない。\par

\noindent\textbf{3：不適切な解答候補} N=0／N>0／Nを2で割った商。初期値11が正の段階で終了し、桁を生成できない。\par

\noindent\textbf{4：不適切な解答候補} N=0／N>0／Nを2で割った余り。分岐が逆であり、更新も桁を消去する商ではない。\par

\noindent\textbf{5：不適切な解答候補} N>0／N=0／2N。Nを増加させるため正の入力で終了に近づかない。\par

\paragraph{再利用する考え方} 位取りの基数と、小数の有限表現が可能かどうかを区別する。小数は2倍操作で上位桁から求める。

\paragraph{共通理論} \hyperref[theory:information:radix]{第
ef*{theory:information:radix}節}

\par 難易度 2（校準初稿）、優先度未評価。

\subsection{平成29年度 問2-2}
\label{q:h29-reg-basic-2-2}

\noindent\textbf{原典：} sources/H29/basic.pdf、PDF物理7頁。

\textbf{公式正答：} 4（single）。

\paragraph{問題の要旨} 13.0を、符号1ビット・指数8ビット・仮数23ビットの指定浮動小数点形式へ変換する。

\paragraph{問われる核心} 正規化と指数バイアス

\paragraph{必要知識} 単精度の3欄、指数バイアス127、正規化先頭1の省略。

\paragraph{考え方（査読待ち初稿）} \(13=8(1+0.625)=2^3(1+0.101_2)\)である。正なので符号0、指数欄は3+127=130=(10000010)₂。仮数欄は先頭の暗黙の1を保存せず、0.625の小数桁101に0を20個加えた23ビットとなる。

\paragraph{正答を選ぶ理由} \(13=8(1+0.625)=2^3(1+0.101_2)\)である。正なので符号0、指数欄は3+127=130=(10000010)₂。仮数欄は先頭の暗黙の1を保存せず、0.625の小数桁101に0を20個加えた23ビットとなる。

\paragraph{選択肢の検討（査読待ち初稿）}

\noindent\textbf{1：不適切な解答候補} 1／10000001／10010000000000000000000。13.0は正なので符号は0である。指数3へ127を加えると130で、10000010になる。0.625=(0.101)₂なので仮数部は101の後に0が20個である。\par

\noindent\textbf{2：不適切な解答候補} 1／10000010／10100000000000000000000。13.0は正なので符号は0である。\par

\noindent\textbf{3：不適切な解答候補} 0／10000001／10100000000000000000000。指数3へ127を加えると130で、10000010になる。\par

\noindent\textbf{4：適切な解答候補} 0／10000010／10100000000000000000000。符号・指数・仮数の3欄がすべて一致する。\par

\noindent\textbf{5：不適切な解答候補} 0／10000001／10100000000000000000000。指数3へ127を加えると130で、10000010になる。\par

\paragraph{再利用する考え方} 先頭の1を除く仮数部と、バイアスを加えた指数部を別々に作る。符号を最後まで追跡する。

\paragraph{共通理論} \hyperref[theory:information:floating]{第
ef*{theory:information:floating}節}

\par 参考：\cite{oracle_single}。

\par 難易度 2（校準初稿）、優先度未評価。

\subsection{平成30年度 問2-3}
\label{q:h30-reg-basic-2-3}

\noindent\textbf{原典：} sources/H30/basic.pdf、PDF物理8頁。

\textbf{公式正答：} 3（single）。

\paragraph{問題の要旨} 01011の5桁を固定して、2の補数と1の補数を求める。

\paragraph{問われる核心} 桁数を固定した補数と負数表現

\paragraph{必要知識} 基数の補数と基数から1を引いた補数、固定桁幅。

\paragraph{考え方（査読待ち初稿）} 01011の十進値は11である。2の補数は32-11=21、10101となる。1の補数は31-11=20、10100となる。後者は各ビットを反転し、前者はそこへ1を加えても得られる。

\paragraph{正答を選ぶ理由} 01011の十進値は11である。2の補数は32-11=21、10101となる。1の補数は31-11=20、10100となる。後者は各ビットを反転し、前者はそこへ1を加えても得られる。

\paragraph{選択肢の検討（査読待ち初稿）}

\noindent\textbf{1：不適切な解答候補} 11011／10100。2の補数は10101であり、11011ではない。\par

\noindent\textbf{2：不適切な解答候補} 10101／11011。1の補数は10100であり、11011ではない。\par

\noindent\textbf{3：適切な解答候補} 10101／10100。2の補数と1の補数がともに一致する。\par

\noindent\textbf{4：不適切な解答候補} 10100／10101。2の補数は10101であり、10100ではない。1の補数は10100であり、10101ではない。\par

\noindent\textbf{5：不適切な解答候補} 11011／11011。2の補数は10101であり、11011ではない。1の補数は10100であり、11011ではない。\par

\paragraph{再利用する考え方} 桁数を固定してから補数を計算する。1の補数はビット反転、2の補数はさらに1を加える。

\paragraph{共通理論} \hyperref[theory:information:complement]{第
ef*{theory:information:complement}節}

\par 難易度 2（校準初稿）、優先度未評価。

\subsection{令和元年度 問2-1}
\label{q:r01-reg-basic-2-1}

\noindent\textbf{原典：} sources/R01/basic.pdf、PDF物理7頁。

\textbf{公式正答：} 1（single）。

\paragraph{問題の要旨} 10進11.5を二進数と16進数へ変換する。

\paragraph{問われる核心} 二進と16進の位取り対応

\paragraph{必要知識} 基数の位取り、4ビットと16進1桁の対応。

\paragraph{考え方（査読待ち初稿）} 11=2\textsuperscript{3}+2+1より整数部1011、0.5=2\textsuperscript{-1}より小数部.1となる。16進整数11はBで、小数0.5を16倍すると8なのでB.8である。二進小数.1を.1000と4ビットにそろえて8へ変換してもよい。

\paragraph{正答を選ぶ理由} 11=2\textsuperscript{3}+2+1より整数部1011、0.5=2\textsuperscript{-1}より小数部.1となる。16進整数11はBで、小数0.5を16倍すると8なのでB.8である。二進小数.1を.1000と4ビットにそろえて8へ変換してもよい。

\paragraph{選択肢の検討（査読待ち初稿）}

\noindent\textbf{1：適切な解答候補} 1011.1／B.8。11=8+2+1、0.5=1/2であり、16進では11=B、0.5=8/16になる。\par

\noindent\textbf{2：不適切な解答候補} 1011.0／C.8。整数12を表すCと、小数0を表す.0が誤りである。\par

\noindent\textbf{3：不適切な解答候補} 1011.1／B.5。16進.5は5/16であり、十進0.5とは異なる。\par

\noindent\textbf{4：不適切な解答候補} 1011.0／B.8。二進.0は小数0であり、0.5を失っている。\par

\noindent\textbf{5：不適切な解答候補} 1011.1／C.5。整数C=12、小数.5=5/16の両方が異なる。\par

\paragraph{再利用する考え方} 位取りの基数と、小数の有限表現が可能かどうかを区別する。小数は2倍操作で上位桁から求める。

\paragraph{共通理論} \hyperref[theory:information:radix]{第
ef*{theory:information:radix}節}

\par 難易度 2（校準初稿）、優先度未評価。

\subsection{令和元年度再試験 問2-3}
\label{q:r01-ret-basic-2-3}

\noindent\textbf{原典：} sources/R01\_retest/basic.pdf、PDF物理8頁。

\textbf{公式正答：} 2（single）。

\paragraph{問題の要旨} 10進表記で2 TBの記憶容量を、問題が定義する2進表記のTBへ換算する。

\paragraph{問われる核心} 情報容量の接頭語と単位換算

\paragraph{必要知識} TBとTiBの単位量、情報量保存。

\paragraph{考え方（査読待ち初稿）} 容量をバイトへ統一してから新単位で割る。\(2\times10^{12}/2^{40}=2(1000/1024)^4\simeq1.819\)であり、約1.8となる。原典の「2進のTB」を現在の用語で明確に区別すればTiBである。

\paragraph{正答を選ぶ理由} 容量をバイトへ統一してから新単位で割る。\(2\times10^{12}/2^{40}=2(1000/1024)^4\simeq1.819\)であり、約1.8となる。原典の「2進のTB」を現在の用語で明確に区別すればTiBである。

\paragraph{選択肢の検討（査読待ち初稿）}

\noindent\textbf{1：不適切な解答候補} 1.6 TB。1.6 TBは2×10\textasciicircum{}12/2\textasciicircum{}40≈1.819の概算と合わない。単位当たりの容量が大きくなるので数値は2より小さくなる。\par

\noindent\textbf{2：適切な解答候補} 1.8 TB。2×10\textasciicircum{}12/2\textasciicircum{}40≈1.819を小数1桁へ概算した値に一致する。\par

\noindent\textbf{3：不適切な解答候補} 2.0 TB。2.0 TBは2×10\textasciicircum{}12/2\textasciicircum{}40≈1.819の概算と合わない。単位当たりの容量が大きくなるので数値は2より小さくなる。\par

\noindent\textbf{4：不適切な解答候補} 2.2 TB。2.2 TBは2×10\textasciicircum{}12/2\textasciicircum{}40≈1.819の概算と合わない。単位当たりの容量が大きくなるので数値は2より小さくなる。\par

\noindent\textbf{5：不適切な解答候補} 2.4 TB。2.4 TBは2×10\textasciicircum{}12/2\textasciicircum{}40≈1.819の概算と合わない。単位当たりの容量が大きくなるので数値は2より小さくなる。\par

\paragraph{再利用する考え方} 単位あたりのバイト数を先に定義し、容量の数値だけを比べない。

\paragraph{共通理論} \hyperref[theory:information:capacity]{第
ef*{theory:information:capacity}節}

\par 参考：\cite{nist_binary_prefix}。

\par 難易度 2（校準初稿）、優先度未評価。

\subsection{令和元年度再試験 問2-4}
\label{q:r01-ret-basic-2-4}

\noindent\textbf{原典：} sources/R01\_retest/basic.pdf、PDF物理9頁。

\textbf{公式正答：} 5（single）。

\paragraph{問題の要旨} 13.0を、符号1ビット・指数8ビット・仮数23ビットの指定浮動小数点形式へ変換する。

\paragraph{問われる核心} 正規化と指数バイアス

\paragraph{必要知識} 単精度の3欄、指数バイアス127、正規化先頭1の省略。

\paragraph{考え方（査読待ち初稿）} \(13=8(1+0.625)=2^3(1+0.101_2)\)である。正なので符号0、指数欄は3+127=130=(10000010)₂。仮数欄は先頭の暗黙の1を保存せず、0.625の小数桁101に0を20個加えた23ビットとなる。

\paragraph{正答を選ぶ理由} \(13=8(1+0.625)=2^3(1+0.101_2)\)である。正なので符号0、指数欄は3+127=130=(10000010)₂。仮数欄は先頭の暗黙の1を保存せず、0.625の小数桁101に0を20個加えた23ビットとなる。

\paragraph{選択肢の検討（査読待ち初稿）}

\noindent\textbf{1：不適切な解答候補} 1／10000010／10100000000000000000000。13.0は正なので符号は0である。\par

\noindent\textbf{2：不適切な解答候補} 1／10000001／10100000000000000000000。13.0は正なので符号は0である。指数3へ127を加えると130で、10000010になる。\par

\noindent\textbf{3：不適切な解答候補} 0／10000001／10100000000000000000000。指数3へ127を加えると130で、10000010になる。\par

\noindent\textbf{4：不適切な解答候補} 0／10000001／10100000000000000000000。指数3へ127を加えると130で、10000010になる。\par

\noindent\textbf{5：適切な解答候補} 0／10000010／10100000000000000000000。符号・指数・仮数の3欄がすべて一致する。\par

\paragraph{再利用する考え方} 先頭の1を除く仮数部と、バイアスを加えた指数部を別々に作る。符号を最後まで追跡する。

\paragraph{共通理論} \hyperref[theory:information:floating]{第
ef*{theory:information:floating}節}

\par 参考：\cite{oracle_single}。

\par 難易度 2（校準初稿）、優先度未評価。

\subsection{令和2年度 問2-4}
\label{q:r02-reg-basic-2-4}

\noindent\textbf{原典：} sources/R02/basic.pdf、PDF物理8頁。

\textbf{公式正答：} 2（single）。

\paragraph{問題の要旨} 100110の6桁を固定して、2の補数と1の補数を求める。

\paragraph{問われる核心} 桁数を固定した補数と負数表現

\paragraph{必要知識} 基数の補数と基数から1を引いた補数、固定桁幅。

\paragraph{考え方（査読待ち初稿）} 100110の十進値は38である。2の補数は64-38=26、011010となる。1の補数は63-38=25、011001となる。後者は各ビットを反転し、前者はそこへ1を加えても得られる。

\paragraph{正答を選ぶ理由} 100110の十進値は38である。2の補数は64-38=26、011010となる。1の補数は63-38=25、011001となる。後者は各ビットを反転し、前者はそこへ1を加えても得られる。

\paragraph{選択肢の検討（査読待ち初稿）}

\noindent\textbf{1：不適切な解答候補} 000110／000101。2の補数は011010であり、000110ではない。1の補数は011001であり、000101ではない。\par

\noindent\textbf{2：適切な解答候補} 011010／011001。2の補数と1の補数がともに一致する。\par

\noindent\textbf{3：不適切な解答候補} 000111／000110。2の補数は011010であり、000111ではない。1の補数は011001であり、000110ではない。\par

\noindent\textbf{4：不適切な解答候補} 011001／011010。2の補数は011010であり、011001ではない。1の補数は011001であり、011010ではない。\par

\noindent\textbf{5：不適切な解答候補} 011000／011001。2の補数は011010であり、011000ではない。\par

\paragraph{再利用する考え方} 桁数を固定してから補数を計算する。1の補数はビット反転、2の補数はさらに1を加える。

\paragraph{共通理論} \hyperref[theory:information:complement]{第
ef*{theory:information:complement}節}

\par 難易度 2（校準初稿）、優先度未評価。

\subsection{令和2年度 問2-5}
\label{q:r02-reg-basic-2-5}

\noindent\textbf{原典：} sources/R02/basic.pdf、PDF物理9頁。

\textbf{公式正答：} 5（single）。

\paragraph{問題の要旨} 11010101を上位桁から二進→十進変換する流れ図の途中値を求める。

\paragraph{問われる核心} 桁評価の漸化式とループ不変条件

\paragraph{必要知識} 位置記数法、s←2s+a\_i、上位からの評価、添字の下降。

\paragraph{考え方（査読待ち初稿）} 上位から1ビット読むたび、既読の値を2倍して新しいビットを加える。更新列は1→3→6→13→26→53→106→213であり、空欄は26／53／106となる。最後だけでなく全段階の添字とビットを対応させる。

\paragraph{正答を選ぶ理由} 上位から1ビット読むたび、既読の値を2倍して新しいビットを加える。更新列は1→3→6→13→26→53→106→213であり、空欄は26／53／106となる。最後だけでなく全段階の添字とビットを対応させる。

\paragraph{選択肢の検討（査読待ち初稿）}

\noindent\textbf{1：不適切な解答候補} 25／52／105。アは26であり25ではない。イは53であり52ではない。ウは106であり105ではない。\par

\noindent\textbf{2：不適切な解答候補} 25／52／106。アは26であり25ではない。イは53であり52ではない。\par

\noindent\textbf{3：不適切な解答候補} 26／52／105。イは53であり52ではない。ウは106であり105ではない。\par

\noindent\textbf{4：不適切な解答候補} 26／53／105。ウは106であり105ではない。\par

\noindent\textbf{5：適切な解答候補} 26／53／106。全更新値が漸化式s←2s+aᵢに一致する。\par

\paragraph{再利用する考え方} 各更新後の値が読了したビット列の値であるという不変条件を用い、最後の一致だけで途中を省略しない。

\paragraph{共通理論} \hyperref[theory:information:horner]{第
ef*{theory:information:horner}節}

\par 難易度 2（校準初稿）、優先度未評価。

\subsection{令和4年度 問2-5}
\label{q:r04-reg-basic-2-5}

\noindent\textbf{原典：} sources/R04/basic.pdf、PDF物理8頁。

\textbf{公式正答：} 5（single）。

\paragraph{問題の要旨} 11001011を上位桁から二進→十進変換する流れ図の途中値を求める。

\paragraph{問われる核心} 桁評価の漸化式とループ不変条件

\paragraph{必要知識} 位置記数法、s←2s+a\_i、上位からの評価、添字の下降。

\paragraph{考え方（査読待ち初稿）} 上位から1ビット読むたび、既読の値を2倍して新しいビットを加える。更新列は1→3→6→12→25→50→101→203であり、空欄は12／25／50／101となる。最後だけでなく全段階の添字とビットを対応させる。

\paragraph{正答を選ぶ理由} 上位から1ビット読むたび、既読の値を2倍して新しいビットを加える。更新列は1→3→6→12→25→50→101→203であり、空欄は12／25／50／101となる。最後だけでなく全段階の添字とビットを対応させる。

\paragraph{選択肢の検討（査読待ち初稿）}

\noindent\textbf{1：不適切な解答候補} 12／25／51／102。ウは50であり51ではない。エは101であり102ではない。\par

\noindent\textbf{2：不適切な解答候補} 13／26／50／102。アは12であり13ではない。イは25であり26ではない。エは101であり102ではない。\par

\noindent\textbf{3：不適切な解答候補} 13／26／52／101。アは12であり13ではない。イは25であり26ではない。ウは50であり52ではない。\par

\noindent\textbf{4：不適切な解答候補} 13／25／50／101。アは12であり13ではない。\par

\noindent\textbf{5：適切な解答候補} 12／25／50／101。全更新値が漸化式s←2s+aᵢに一致する。\par

\paragraph{再利用する考え方} 各更新後の値が読了したビット列の値であるという不変条件を用い、最後の一致だけで途中を省略しない。

\paragraph{共通理論} \hyperref[theory:information:horner]{第
ef*{theory:information:horner}節}

\par 難易度 2（校準初稿）、優先度未評価。

\subsection{令和6年度 問2-1}
\label{q:r06-reg-basic-2-1}

\noindent\textbf{原典：} sources/R06/basic.pdf、PDF物理6頁。

\textbf{公式正答：} 1（single）。

\paragraph{問題の要旨} 10進0.6を二進小数で小数点以下16桁まで表す。

\paragraph{問われる核心} 循環二進小数と表示精度

\paragraph{必要知識} 反復2倍法、剰余の循環。

\paragraph{考え方（査読待ち初稿）} 0.6を2倍して整数1、残り0.2から0、残り0.4から0、残り0.8から1、残り0.6へ戻る。したがって1001を反復し、小数16桁は1001100110011001となる。分母5は2の累乗ではないので展開は有限ではない。

\paragraph{正答を選ぶ理由} 0.6を2倍して整数1、残り0.2から0、残り0.4から0、残り0.8から1、残り0.6へ戻る。したがって1001を反復し、小数16桁は1001100110011001となる。分母5は2の累乗ではないので展開は有限ではない。

\paragraph{選択肢の検討（査読待ち初稿）}

\noindent\textbf{1：適切な解答候補} 0.1001100110011001。0.6の反復2倍から得る1001100110011001に一致する。\par

\noindent\textbf{2：不適切な解答候補} 0.1011001100110011。最初の3桁を101としているが、0.6では100となる。\par

\noindent\textbf{3：不適切な解答候補} 0.1100000000000000。0.75の有限二進表現であり、0.6と異なる。\par

\noindent\textbf{4：不適切な解答候補} 0.1100110011001100。先頭2桁が11なら0.75以上となり、0.6と合わない。\par

\noindent\textbf{5：不適切な解答候補} 0.1110011001100110。先頭3桁が111なら0.875以上となり、0.6と合わない。\par

\paragraph{再利用する考え方} 位取りの基数と、小数の有限表現が可能かどうかを区別する。小数は2倍操作で上位桁から求める。

\paragraph{共通理論} \hyperref[theory:information:radix]{第
ef*{theory:information:radix}節}

\par 難易度 2（校準初稿）、優先度未評価。

\section{アルゴリズム・再帰}

\subsection{平成24年度 問2-5}
\label{q:h24-reg-basic-2-5}

\noindent\textbf{原典：} sources/H24/basic.pdf、PDF物理7頁。

\textbf{公式正答：} 2（single）。

\par\textit{原典本文の校正と独自解説は未作成。分類のみの登録項目。}

\subsection{平成26年度 問2-2}
\label{q:h26-reg-basic-2-2}

\noindent\textbf{原典：} sources/H26/basic.pdf、PDF物理6頁。

\textbf{公式正答：} 2（single）。

\par\textit{原典本文の校正と独自解説は未作成。分類のみの登録項目。}

\subsection{平成29年度 問2-3}
\label{q:h29-reg-basic-2-3}

\noindent\textbf{原典：} sources/H29/basic.pdf、PDF物理8頁。

\textbf{公式正答：} 5（single）。

\par\textit{原典本文の校正と独自解説は未作成。分類のみの登録項目。}

\subsection{令和元年度再試験 問2-2}
\label{q:r01-ret-basic-2-2}

\noindent\textbf{原典：} sources/R01\_retest/basic.pdf、PDF物理7頁。

\textbf{公式正答：} 4（single）。

\par\textit{原典本文の校正と独自解説は未作成。分類のみの登録項目。}

\subsection{令和3年度 問2-6}
\label{q:r03-reg-basic-2-6}

\noindent\textbf{原典：} sources/R03/basic.pdf、PDF物理10頁。

\textbf{公式正答：} 3（single）。

\par\textit{原典本文の校正と独自解説は未作成。分類のみの登録項目。}

\subsection{令和4年度 問2-3}
\label{q:r04-reg-basic-2-3}

\noindent\textbf{原典：} sources/R04/basic.pdf、PDF物理6頁。

\textbf{公式正答：} 3（single）。

\par\textit{原典本文の校正と独自解説は未作成。分類のみの登録項目。}

\subsection{令和5年度 問2-2}
\label{q:r05-reg-basic-2-2}

\noindent\textbf{原典：} sources/R05/basic.pdf、PDF物理6頁。

\textbf{公式正答：} 4（single）。

\par\textit{原典本文の校正と独自解説は未作成。分類のみの登録項目。}

\subsection{令和6年度 問2-5}
\label{q:r06-reg-basic-2-5}

\noindent\textbf{原典：} sources/R06/basic.pdf、PDF物理8頁。

\textbf{公式正答：} 4（single）。

\par\textit{原典本文の校正と独自解説は未作成。分類のみの登録項目。}

\subsection{令和7年度 問2-6}
\label{q:r07-reg-basic-2-6}

\noindent\textbf{原典：} sources/R07/basic.pdf、PDF物理7頁。

\textbf{公式正答：} 1（single）。

\par\textit{原典本文の校正と独自解説は未作成。分類のみの登録項目。}

\section{形式言語・木}

\subsection{平成24年度 問2-2}
\label{q:h24-reg-basic-2-2}

\noindent\textbf{原典：} sources/H24/basic.pdf、PDF物理5頁。

\textbf{公式正答：} 2（single）。

\par\textit{原典本文の校正と独自解説は未作成。分類のみの登録項目。}

\subsection{平成26年度 問2-5}
\label{q:h26-reg-basic-2-5}

\noindent\textbf{原典：} sources/H26/basic.pdf、PDF物理8頁。

\textbf{公式正答：} 5（single）。

\par\textit{原典本文の校正と独自解説は未作成。分類のみの登録項目。}

\subsection{平成27年度 問2-6}
\label{q:h27-reg-basic-2-6}

\noindent\textbf{原典：} sources/H27/basic.pdf、PDF物理7頁。

\textbf{公式正答：} 4（single）。

\par\textit{原典本文の校正と独自解説は未作成。分類のみの登録項目。}

\subsection{平成29年度 問2-5}
\label{q:h29-reg-basic-2-5}

\noindent\textbf{原典：} sources/H29/basic.pdf、PDF物理10頁。

\textbf{公式正答：} 4（single）。

\par\textit{原典本文の校正と独自解説は未作成。分類のみの登録項目。}

\subsection{令和3年度 問2-5}
\label{q:r03-reg-basic-2-5}

\noindent\textbf{原典：} sources/R03/basic.pdf、PDF物理9頁。

\textbf{公式正答：} 1（single）。

\par\textit{原典本文の校正と独自解説は未作成。分類のみの登録項目。}

\section{ベクトル・微分演算}

\subsection{令和元年度 問2-3}
\label{q:r01-reg-basic-2-3}

\noindent\textbf{原典：} sources/R01/basic.pdf、PDF物理9頁。

\textbf{公式正答：} 3（single）。

\par\textit{原典本文の校正と独自解説は未作成。分類のみの登録項目。}

\section{環境法・循環・持続可能性}

\subsection{令和4年度 問2-1}
\label{q:r04-reg-basic-2-1}

\noindent\textbf{原典：} sources/R04/basic.pdf、PDF物理5頁。

\textbf{公式正答：} 1（single）。

\par\textit{原典本文の校正と独自解説は未作成。分類のみの登録項目。}

\section{情報理論・符号・誤り検出}

\subsection{平成26年度 問2-3}
\label{q:h26-reg-basic-2-3}

\noindent\textbf{原典：} sources/H26/basic.pdf、PDF物理7頁。

\textbf{公式正答：} 4（single）。

\par\textit{原典本文の校正と独自解説は未作成。分類のみの登録項目。}

\subsection{平成28年度 問2-5}
\label{q:h28-reg-basic-2-5}

\noindent\textbf{原典：} sources/H28/basic.pdf、PDF物理10頁。

\textbf{公式正答：} 4（single）。

\par\textit{原典本文の校正と独自解説は未作成。分類のみの登録項目。}

\subsection{令和5年度 問2-3}
\label{q:r05-reg-basic-2-3}

\noindent\textbf{原典：} sources/R05/basic.pdf、PDF物理7頁。

\textbf{公式正答：} 5（single）。

\par\textit{原典本文の校正と独自解説は未作成。分類のみの登録項目。}

\subsection{令和6年度 問2-4}
\label{q:r06-reg-basic-2-4}

\noindent\textbf{原典：} sources/R06/basic.pdf、PDF物理7頁。

\textbf{公式正答：} 5（single）。

\par\textit{原典本文の校正と独自解説は未作成。分類のみの登録項目。}

\subsection{令和7年度 問2-2}
\label{q:r07-reg-basic-2-2}

\noindent\textbf{原典：} sources/R07/basic.pdf、PDF物理5頁。

\textbf{公式正答：} 4（single）。

\par\textit{原典本文の校正と独自解説は未作成。分類のみの登録項目。}

\section{スタック・データ構造}

\subsection{平成26年度 問2-4}
\label{q:h26-reg-basic-2-4}

\noindent\textbf{原典：} sources/H26/basic.pdf、PDF物理7頁。

\textbf{公式正答：} 4（single）。

\paragraph{問題の要旨} 空スタックへ指定された8操作を実行し、最後のPOPの値を求める。

\paragraph{問われる核心} データ構造・形式言語

\paragraph{必要知識} LIFO、PUSH/POP、状態追跡

\paragraph{考え方（査読待ち初稿）} スタックを底から上へ記録する。操作後は [1]→[1,2]→[1,2,3]→[1,2]→[1,2,4]→[1,2,4,5]→[1,2,4]→[1,2]。最後のPOPで取り出す値は4なので選択肢④。

\paragraph{他の選択肢（査読待ち初稿）} ①1と②2は最後のPOP後も残る。③3は4番目の操作で、⑤5は7番目の操作で取り出される。途中のPOPと最後のPOPを混同しない。

\par 難易度 2（校準初稿）、優先度未評価。

\subsection{令和元年度 問2-6}
\label{q:r01-reg-basic-2-6}

\noindent\textbf{原典：} sources/R01/basic.pdf、PDF物理11頁。

\textbf{公式正答：} 2（single）。

\par\textit{原典本文の校正と独自解説は未作成。分類のみの登録項目。}

\section{未確定}

\subsection{平成24年度 問2-3}
\label{q:h24-reg-basic-2-3}

\noindent\textbf{原典：} sources/H24/basic.pdf、PDF物理6頁。

\textbf{公式正答：} 4（single）。

\par\textit{原典本文の校正と独自解説は未作成。分類のみの登録項目。}

\subsection{平成24年度 問2-4}
\label{q:h24-reg-basic-2-4}

\noindent\textbf{原典：} sources/H24/basic.pdf、PDF物理6頁。

\textbf{公式正答：} 3（single）。

\par\textit{原典本文の校正と独自解説は未作成。分類のみの登録項目。}

\subsection{平成25年度 問2-1}
\label{q:h25-reg-basic-2-1}

\noindent\textbf{原典：} sources/H25/basic.pdf、PDF物理6頁。

\textbf{公式正答：} 4（single）。

\par\textit{原典本文の校正と独自解説は未作成。分類のみの登録項目。}

\subsection{平成25年度 問2-2}
\label{q:h25-reg-basic-2-2}

\noindent\textbf{原典：} sources/H25/basic.pdf、PDF物理7頁。

\textbf{公式正答：} 5（single）。

\par\textit{原典本文の校正と独自解説は未作成。分類のみの登録項目。}

\subsection{平成27年度 問2-2}
\label{q:h27-reg-basic-2-2}

\noindent\textbf{原典：} sources/H27/basic.pdf、PDF物理5頁。

\textbf{公式正答：} 5（single）。

\par\textit{原典本文の校正と独自解説は未作成。分類のみの登録項目。}

\subsection{平成28年度 問2-1}
\label{q:h28-reg-basic-2-1}

\noindent\textbf{原典：} sources/H28/basic.pdf、PDF物理7頁。

\textbf{公式正答：} 4（single）。

\par\textit{原典本文の校正と独自解説は未作成。分類のみの登録項目。}

\subsection{平成30年度 問2-5}
\label{q:h30-reg-basic-2-5}

\noindent\textbf{原典：} sources/H30/basic.pdf、PDF物理9頁。

\textbf{公式正答：} 1（single）。

\par\textit{原典本文の校正と独自解説は未作成。分類のみの登録項目。}

\subsection{平成30年度 問2-6}
\label{q:h30-reg-basic-2-6}

\noindent\textbf{原典：} sources/H30/basic.pdf、PDF物理10頁。

\textbf{公式正答：} 2（single）。

\par\textit{原典本文の校正と独自解説は未作成。分類のみの登録項目。}

\subsection{令和元年度 問2-2}
\label{q:r01-reg-basic-2-2}

\noindent\textbf{原典：} sources/R01/basic.pdf、PDF物理8頁。

\textbf{公式正答：} 1（single）。

\par\textit{原典本文の校正と独自解説は未作成。分類のみの登録項目。}

\subsection{令和元年度 問2-4}
\label{q:r01-reg-basic-2-4}

\noindent\textbf{原典：} sources/R01/basic.pdf、PDF物理9頁。

\textbf{公式正答：} 3（single）。

\par\textit{原典本文の校正と独自解説は未作成。分類のみの登録項目。}

\subsection{令和元年度再試験 問2-6}
\label{q:r01-ret-basic-2-6}

\noindent\textbf{原典：} sources/R01\_retest/basic.pdf、PDF物理10頁。

\textbf{公式正答：} 4（single）。

\par\textit{原典本文の校正と独自解説は未作成。分類のみの登録項目。}

\subsection{令和2年度 問2-1}
\label{q:r02-reg-basic-2-1}

\noindent\textbf{原典：} sources/R02/basic.pdf、PDF物理7頁。

\textbf{公式正答：} 4（single）。

\par\textit{原典本文の校正と独自解説は未作成。分類のみの登録項目。}

\subsection{令和4年度 問2-2}
\label{q:r04-reg-basic-2-2}

\noindent\textbf{原典：} sources/R04/basic.pdf、PDF物理5頁。

\textbf{公式正答：} 3（single）。

\par\textit{原典本文の校正と独自解説は未作成。分類のみの登録項目。}

\subsection{令和5年度 問2-1}
\label{q:r05-reg-basic-2-1}

\noindent\textbf{原典：} sources/R05/basic.pdf、PDF物理5頁。

\textbf{公式正答：} 1（single）。

\par\textit{原典本文の校正と独自解説は未作成。分類のみの登録項目。}

\subsection{令和5年度 問2-4}
\label{q:r05-reg-basic-2-4}

\noindent\textbf{原典：} sources/R05/basic.pdf、PDF物理7頁。

\textbf{公式正答：} 1（single）。

\par\textit{原典本文の校正と独自解説は未作成。分類のみの登録項目。}

\subsection{令和5年度 問2-6}
\label{q:r05-reg-basic-2-6}

\noindent\textbf{原典：} sources/R05/basic.pdf、PDF物理8頁。

\textbf{公式正答：} 4（single）。

\par\textit{原典本文の校正と独自解説は未作成。分類のみの登録項目。}

\subsection{令和6年度 問2-2}
\label{q:r06-reg-basic-2-2}

\noindent\textbf{原典：} sources/R06/basic.pdf、PDF物理6頁。

\textbf{公式正答：} 2（single）。

\par\textit{原典本文の校正と独自解説は未作成。分類のみの登録項目。}

\subsection{令和7年度 問2-1}
\label{q:r07-reg-basic-2-1}

\noindent\textbf{原典：} sources/R07/basic.pdf、PDF物理5頁。

\textbf{公式正答：} 3（single）。

\par\textit{原典本文の校正と独自解説は未作成。分類のみの登録項目。}
